\section{Motivation}
% \Khang{Nếu đổi thành đề tài Information Retrieval $\rightarrow$ Sửa lại cho phù hợp}

% \Khang{Cần thiết fashion nữa không? Có thể cần nhưng phải chú trọng feature extraction và multimodal attention}
% % In the evolving landscape of economics and retail, the fashion industry is experiencing rapid growth, catering to diverse demographics and preferences. However, as the number of fashion stores increases, consumers often find themselves spending considerable time searching for a store that aligns with their unique interests, financial constraints, and other specific requirements. This search predominantly relies on human-assisted interactions through websites, where customers engage in conversations with store managers.  To address this challenge, the development of a multimodal chatbot emerges as a potential solution. By integrating image and text into the conversation, the chatbot can efficiently assist customers in navigating through the diverse offerings and overall product information of the fashion stores. 

% \noindent As the fashion market expands to cater to diverse demographics and preferences, the conventional methods of searching for products become increasingly inefficient. Consumers often face challenges in finding fashion items that align with their specific preferences, budget constraints, and other unique requirements. Traditional approaches rely heavily on human-assisted interactions through websites, where customers engage in conversations with store managers to seek guidance or recommendations. However, this manual process is time-consuming, prone to errors, and fails to leverage the full potential of technology in enhancing customer experiences. The inefficiency of this process raises a significant challenge: How can this whole process be automated to enhance customer experience and streamline operations for fashion stores?

% \noindent By introducing a multimodal searching method, we aim to revolutionize the way customers interact with fashion stores online. The integration of image and text capabilities allows for a more intuitive and efficient communication channel between the consumer and the store. With this technology, customers can simply provide an image or description of the desired product, and the system can quickly analyze their preferences and recommend suitable options. Through fine-tuning a Vietnamese multimodal model and implementing an efficient information retrieval pipeline, the system can understand customer preferences, recommend appropriate products, and provide personalized assistance, thus revolutionizing the way customers interact with fashion stores online.

% \noindent Moreover, the development of such a system addresses the pressing need for automation in the fashion retail sector. As the number of fashion stores continues to rise, streamlining operations and improving customer experiences become paramount. A well-trained multimodal search system can effectively assist customers in navigating through the vast array of products, providing personalized recommendations, and answering queries regarding price, availability, and discounts.


